\section{Methods}
\label{sec:methods}

\subsection{Architecture}
\label{sec:architecture}

Game states are graphs with one node per crossing. Node features carry the
ordered PD tuple plus the ternary resolution state; edges pair the two
occurrences of each arc label with the strand identity and its under/over
resolution status. The primary model (\texttt{PortGraphTransformerNet} in
\texttt{knot\_graph\_nnet.py}) first runs typed message passing over the four
half-edges (ports) of every crossing---self, cyclic neighbors, opposite
strand, and same-arc peer---then pools ports to crossing tokens, adds a direct
crossing-state residual and a positional signal, and applies a global
transformer with a graph token for the value head and per-crossing keep/switch
logits for the policy head. Positions use deterministic first-encounter PD
traversal ranks (or, in the variable-size model, scale-free sinusoidal/linear
traversal features with no crossing-count-indexed lookup), so identity derives
from diagram topology rather than arbitrary list order. Fixed-shadow controls
include a crossing-state MLP and a dense PD-state MLP
(\texttt{capacity\_test.py}); the AlphaZero loop itself follows the standard
pattern---MCTS with network priors, self-play episode generation, network
training, and arena gating of champions (\texttt{mcts.py}, \texttt{coach.py},
\texttt{arena.py}).

\subsection{Corpus and manifest}
\label{sec:corpus}

\texttt{prime\_knot\_corpus.py} stores 35 standard Spherogram table diagrams,
one per prime knot type from $3_1$ through $8_{21}$ ($3{:}1$, $4{:}1$, $5{:}2$,
$6{:}3$, $7{:}7$, $8{:}21$). These are table \emph{diagrams}, not an
enumeration of shadows; strategy formally belongs to a shadow game, and any
knot-type-level claim requires robustness across diagrams.
\texttt{corpus\_manifest.py} assigns each record a stable \texttt{shadow\_id}
(hash of the canonical PD), source identifier, canonicalization version,
equivalence group over relabelings and permutations, frozen split, and
exact-table availability with an on-demand exact-table checksum
(\texttt{prime-catalog-v1}). Splits are knot-type-disjoint with singletons kept
in train so every crossing count is represented; canonical relabelings of one
object can never leak across a split (\texttt{assert\_no\_leakage}).

\subsection{Training regimes}
\label{sec:training}

We distinguish three settings (ROADMAP G1.3): (i) \emph{per-shadow capacity},
exact-table supervision and exhaustive evaluation on one shadow;
(ii) \emph{shared supervised capacity}, one variable-size model fit to exact
tables from many shadows with hierarchical sampling over crossing number, knot
identity, and depth (\texttt{variable\_size\_capacity\_test.py}, supporting
\texttt{--split} for frozen train/validation/test reporting); and
(iii) \emph{held-out generalization}, evaluation on frozen shadows or knot
types whose exact tables supplied no training label. Only (iii) may be called
generalization. Self-play (G1.4) trains without exact targets and is compared
to minimax by raw-network policy, fixed-budget MCTS policy, and exact tables
on every tractable shadow.
